% Relatório: quando orquestrar vários agentes compensa (e quando não), com a literatura e o braço
% "sozinho" do benchmark do NarraTer.
% Compilar: cd docs/relatorios && pdflatex quando-orquestrar.tex && pdflatex quando-orquestrar.tex
% Fontes: docs/pesquisa/orquestracao.md (estudo) e bench/orquestracao/resultados/ (medições).
\documentclass[11pt,a4paper]{article}
\usepackage[utf8]{inputenc}
\usepackage[T1]{fontenc}
\usepackage[brazil]{babel}
\usepackage{lmodern}
\usepackage[margin=2.2cm]{geometry}
\usepackage{booktabs}
\usepackage{tabularx}
\usepackage{xcolor}
\usepackage{graphicx}
\usepackage{float}
\usepackage{enumitem}
\usepackage{caption}
\usepackage[hidelinks]{hyperref}
\setlist{itemsep=2pt,topsep=4pt}
\captionsetup{font=small,labelfont=bf}

\newcommand{\rot}[1]{\textsf{\footnotesize[#1]}}
\newcolumntype{Y}{>{\raggedright\arraybackslash}X}

\title{Quando orquestrar agentes não compensa:\\[2pt]
\large o que a pesquisa mostra e o que medimos no NarraTer}
\author{Gabriel Cortez \and Claude Code}
\date{16 de outubro de 2026}

\begin{document}
\maketitle

\begin{abstract}
\noindent
O NarraTer existe para pôr vários agentes de IA para trabalhar juntos. Este documento faz a
pergunta incômoda: na maior parte das tarefas de programação, isso compensa? Reunimos o que
artigos e equipes de produção mediram em 2025 e 2026 e o comparamos com o nosso próprio benchmark.
A resposta da literatura é consistente. Quando um agente sozinho já resolve bem a tarefa, uma equipe
custa de 1,8 a 5,1 vezes mais e quase nunca melhora o resultado de forma significativa. Em tarefas
sequenciais, ela piora. A equipe compensa em casos estreitos: trabalho realmente paralelo, contexto
que não cabe num agente só e verificação independente. Medimos isso no NarraTer com o mesmo pedido
dado a um terminal sozinho e à equipe de três agentes (seção~\ref{sec:medido}). Os dois tiraram
31 de 31 nos testes escondidos. O terminal sozinho foi 2,5 vezes mais rápido (1,6 contra 3,9~min),
usou 7 vezes menos tokens e custou 26\% menos que a versão atual da equipe (G).
\end{abstract}

\noindent\textbf{Rótulos.} \rot{citado}: afirmado pela fonte indicada. \rot{medido}: medido por nós,
com método reproduzível. \rot{inferido}: nossa leitura, sem medição direta.

\section{A pergunta}

Montar uma equipe de agentes parece óbvio: um líder divide a tarefa, especialistas executam em paralelo
e o resultado sai mais rápido e melhor. Nosso benchmark de versões (relatório
\emph{orquestracao-versoes.pdf}) mediu o \emph{custo} de orquestrar: tokens, tempo e falhas do
protocolo. Ele não mediu se orquestrar \emph{vale a pena}. Todas as rodadas usaram a equipe, e
nenhuma o agente sozinho. Este documento responde a essa segunda pergunta, primeiro com a literatura
e depois com uma medição nossa.

\section{O que a pesquisa mediu}

\subsection{Com o mesmo orçamento, um agente sozinho empata ou ganha}

A comparação justa é a de mesmo orçamento: a equipe gasta mais tokens, então precisa ser comparada
com um agente sozinho que também pode gastar mais (pensar mais, tentar de novo).

\begin{itemize}
  \item \rot{citado} Kim et al.~\cite{kim2025} avaliaram 180 configurações de sistemas de agentes.
    A coordenação dá ``retorno decrescente ou negativo quando o agente sozinho passa de ${\sim}$45\%''
    de acerto. Em tarefas paralelizáveis, a coordenação centralizada ganha 80,8\%. Em raciocínio
    sequencial, as variantes multiagente perdem de 39\% a 70\%.
  \item \rot{citado} Tran e Kiela~\cite{tran2026} mostram que, com o mesmo orçamento de tokens de
    raciocínio, um agente sozinho empata ou supera sistemas multiagente em raciocínio de vários passos.
  \item \rot{citado} Wang et al.~\cite{rethinking2024} chegam ao mesmo resultado: um agente sozinho
    bem instruído chega perto da melhor discussão entre agentes. A discussão só vence quando não há
    exemplos no prompt.
  \item \rot{citado} Zhang et al.~\cite{mad2025} (5 métodos de debate, 9 benchmarks) concluem que o debate entre agentes raramente supera
    o raciocínio encadeado ou a autoconsistência de um agente só. A exceção é usar modelos diferentes
    entre si.
\end{itemize}

\subsection{Em programação, equipes custam mais e quase não melhoram}

A Vals.ai publicou em 9 de outubro de 2026 o estudo mais próximo do nosso caso~\cite{vals2026}.
A tarefa era construir aplicações web inteiras, com nota dada por testes de interface. Cada modelo
rodou sozinho e como equipe (um líder e até cinco subagentes), em dois níveis de esforço.

\begin{table}[H]
\centering\small
\caption{Agente sozinho contra equipe na Vibe Code Bench \rot{citado}~\cite{vals2026}.}
\begin{tabular}{lrrrr}
\toprule
Modelo (esforço) & Sozinho & Equipe & Custo da equipe & Tempo mediano (min) \\
\midrule
GPT 6 Sol (médio)   & 77,6\% & 84,9\% (+7,3*) & 2,5$\times$ & 12,4 $\to$ 11,9 \\
GPT 6 Sol (máximo)  & 89,0\% & 90,4\% (+1,4)  & 1,8$\times$ & 38,4 $\to$ 27,6 \\
Opus 5.5 (médio)    & 91,5\% & 91,2\% ($-$0,3) & 2,2$\times$ & 20,3 $\to$ 26,6 \\
Opus 5.5 (máximo)   & 89,8\% & 93,2\% (+3,4)  & 5,1$\times$ & 74,7 $\to$ 174,1 \\
\bottomrule
\multicolumn{5}{l}{\footnotesize * a única diferença estatisticamente significativa ($p = 0{,}005$).}
\end{tabular}
\end{table}

Das quatro comparações, só uma teve ganho significativo, e foi justamente no modelo com a nota mais
baixa sozinho. Com o Opus 5.5 no esforço médio, a configuração mais parecida com o nosso Narrador, a
equipe ficou 0,3 ponto abaixo, custando 2,2 vezes mais e levando 6 minutos a mais. No esforço máximo,
a equipe do Opus custou US\$~122 por aplicação, contra US\$~23,77 sozinho, e levou 174 minutos
contra 75.

Outras medições em código apontam na mesma direção:
\begin{itemize}
  \item \rot{citado} O Agentless~\cite{agentless2024} é um pipeline fixo (localizar, corrigir,
    validar), sem nenhum agente planejando. Resolveu 32\% do SWE-bench Lite a US\$~0,70 por problema,
    competitivo com sistemas de agentes muito mais caros na época.
  \item \rot{citado} A Anthropic, descrevendo o próprio sistema multiagente de pesquisa~\cite{anthropic2025},
    diz que programação tem ``menos tarefas realmente paralelizáveis que pesquisa'' e que
    multiagente encaixa mal onde todos precisam do mesmo contexto.
  \item \rot{citado} Um estudo de 2026 sobre topologias em código do zero~\cite{topologia2026}
    mostra que só mudar a topologia move a nota em mais de 30 pontos e dobra o tempo. Pipelines
    estruturados convergem mais rápido e com mais qualidade, e ``supervisão gerencial pesada piora o
    desempenho''.
\end{itemize}

\subsection{Por que a equipe custa tanto}

O sistema de pesquisa da Anthropic gasta cerca de 15 vezes os tokens de uma conversa, contra cerca de
4 vezes de um agente sozinho~\cite{anthropic2025} \rot{citado}. Lá, a variação de tokens sozinha explica 80\% da variação
de desempenho. Isso significa que boa parte do ganho de uma equipe é simplesmente gastar mais. Os
mecanismos que a literatura descreve:

\begin{enumerate}
  \item \textbf{Tradução entre agentes.} Cada delegação reescreve a tarefa, e cada resposta é lida e
    resumida de novo pelo líder. No benchmark da LangChain~\cite{langchain2025}, a maior parte dos
    erros do supervisor vinha desse ``telefone sem fio''. Repassar a resposta palavra por palavra
    melhorou o resultado em quase 50\% \rot{citado}.
  \item \textbf{Decisões implícitas em conflito.} A Cognition~\cite{cognition2025} mostra
    subagentes paralelos produzindo peças incompatíveis: ``ações carregam decisões implícitas, e
    decisões em conflito geram resultados ruins''. Em 2026~\cite{cognition2026}, a mesma equipe
    conclui que o multiagente funciona quando ``a escrita fica em uma linha só'' e os outros agentes
    contribuem inteligência, não ações \rot{citado}.
  \item \textbf{Amplificação de erro.} Agentes independentes amplificam erros 17,2 vezes; a
    coordenação centralizada contém isso em 4,4 vezes~\cite{kim2025} \rot{citado}.
  \item \textbf{Verificação fraca.} Na taxonomia MAST~\cite{mast2025}, tirada de 150 traços
    anotados de 7 frameworks, cerca de 44\% das falhas são de projeto do sistema (repetir passos
    15,7\%, não saber quando parar 12,4\%), cerca de 32\% de desalinhamento entre agentes e cerca de
    23,5\% de verificação (verificação errada 9,1\%, ausente 8,2\%, término prematuro 6,2\%). Em
    trajetórias de código, 75,8\% das falhas com declaração de status eram um ``pronto'' falso, e
    juízes LLM não passaram de AUROC 0,65 para detectá-lo~\cite{falsosucesso2026} \rot{citado}.
\end{enumerate}

\subsection{Quando a equipe compensa}

A mesma literatura aponta os casos em que vale montar uma equipe:
\begin{itemize}
  \item \textbf{Trabalho realmente paralelo.} Subtarefas independentes, de preferência de leitura
    (pesquisar, comparar, auditar), com coordenação centralizada: +80,8\%~\cite{kim2025}. A pesquisa
    da Anthropic venceu o Opus sozinho por 90,2\% numa avaliação interna~\cite{anthropic2025} \rot{citado}.
  \item \textbf{Contexto que não cabe num agente.} Com muitas ferramentas ou domínios misturados, um
    agente sozinho ``cai bruscamente'' a partir de dois domínios de distração~\cite{langchain2025} \rot{citado}.
  \item \textbf{Verificação independente.} Um revisor com contexto limpo, que vê só o diff, acha em
    média 2 bugs por PR, cerca de 58\% deles graves~\cite{cognition2026}. Rodar os testes como retorno
    rendeu de 4 a 5 pontos no MetaGPT~\cite{metagpt2023} \rot{citado}.
  \item \textbf{Modelo base fraco na tarefa.} O único ganho significativo da Vals.ai foi no modelo com a
    menor nota sozinho~\cite{vals2026}, coerente com o limiar de ${\sim}$45\% de Kim et al. \rot{inferido}
\end{itemize}

\section{O que medimos no NarraTer}\label{sec:medido}

\subsection{O que o estudo de versões já dizia}

No estudo de versões (A a G, 3 rodadas cada), \textbf{todas} as rodadas tiraram 31 de 31 nos testes
escondidos \rot{medido}. Pelo critério de Kim et al., a tarefa está muito acima do limiar de 45\%: só
custo e tempo podem diferir. O custo cresceu com o número de pedidos que o Narrador fazia aos devs, e a
latência do protocolo em si ficou em segundos. Na melhor versão (G), a equipe gastou 3,9~min e
US\$~0,74 por rodada, em média \rot{medido}.

\subsection{O braço ``sozinho''}

\paragraph{Método.} Usamos o mesmo benchmark, a mesma tarefa (o repositório \texttt{pedidos}, com uma
funcionalidade de várias moedas e dois defeitos plantados) e os mesmos 31 testes escondidos. Cada
rodada começa do zero: HOME, tmux e sessão próprios. Há uma diferença só: o canvas tem \textbf{um
terminal}, com o mesmo Opus 5.5 do Narrador e o papel de desenvolvedor, porque o papel de líder proíbe
editar. O pedido é o mesmo, trocando ``distribua o trabalho entre DevA e DevB'' por ``você trabalha
sozinho''. O critério de pronto e a linha de fim não mudam. Foram 3 rodadas, na versão atual (H,
\texttt{73fab8b}), em 16/10/2026 (\texttt{BENCH\_VARIANTE=solo python3 rodar.py H N}; comparação em
\texttt{bench/orquestracao/solo.py}). Para a equipe, usamos as rodadas F e G do estudo de versões, as
duas melhores. A H só acrescenta o painel Activity à G, sem mudar o protocolo.

\begin{table}[H]
\centering\small
\caption{Um terminal sozinho contra a equipe de três agentes, na mesma tarefa \rot{medido}. Média e,
entre parênteses, o mínimo e o máximo das 3 rodadas. A equipe F tem o mesmo protocolo de leitura do
transcript da G, sem o ajuste de pedidos em lote.}
\begin{tabular}{lrrr}
\toprule
 & Sozinho (H) & Equipe (G) & Equipe (F) \\
\midrule
Testes escondidos (de 31) & 31 (31--31) & 31 (31--31) & 31 (31--31) \\
Duração (min) & 1{,}6 (1{,}4--1{,}8) & 3{,}9 (2{,}8--5{,}2) & 2{,}8 (2{,}3--3{,}7) \\
Tokens processados (milhões) & 0{,}36 (0{,}35--0{,}36) & 2{,}50 (2{,}01--2{,}85) & 1{,}73 (1{,}27--2{,}28) \\
Custo de API equivalente (US\$) & 0{,}55 (0{,}52--0{,}58) & 0{,}74 (0{,}63--0{,}88) & 0{,}65 (0{,}50--0{,}82) \\
Chamadas à API (todos os agentes) & 8 (8--8) & 49 (42--56) & 38 (30--49) \\
Pedidos do Narrador a outros & 0 (0--0) & 4 (3--6) & 4 (3--7) \\
Linhas no diff & 368 (340--412) & 454 (413--478) & 394 (316--524) \\
\bottomrule
\end{tabular}

\end{table}

\paragraph{Resultado.} \rot{medido}
\begin{itemize}
  \item \textbf{Mesma qualidade.} As seis rodadas da equipe e as três rodadas sozinho tiraram 31 de 31.
    Os testes do próprio agente também passaram todos (40 a 47).
  \item \textbf{2,5 vezes mais rápido que a equipe G}: 1,6~min contra 3,9~min, sem sobreposição
    (Mann--Whitney $U = 0$, $p = 0{,}10$, o menor $p$ possível com 3 contra 3). Contra a F, a melhor
    equipe em tempo, o agente sozinho também vence sem sobreposição (1,6 contra 2,8~min, $p = 0{,}10$).
  \item \textbf{7 vezes menos tokens} (0,36 contra 2,50 milhões) e \textbf{6 vezes menos chamadas à
    API} (8 contra 49), também sem sobreposição. As três rodadas sozinho fizeram exatamente 8 chamadas.
  \item \textbf{26\% mais barato que a G} (US\$~0,55 contra 0,74, sem sobreposição, $p = 0{,}10$).
    Contra a F, o custo se sobrepõe (US\$~0,50 a 0,82; $p = 0{,}70$): não dá para afirmar diferença.
\end{itemize}

\paragraph{Por que o custo cai menos que os tokens.} Na equipe, quem processa mais tokens é o DevB,
o Haiku: de 1,3 a 1,7 milhão de tokens por rodada na G, a menos de US\$~0,04. O agente sozinho é o
Opus, o modelo mais caro. Mesmo assim, \textbf{só o Narrador da equipe G, apenas coordenando, custou
de US\$~0,32 a 0,41 por rodada}: de 60\% a 75\% do que o Opus sozinho gastou para fazer a tarefa
inteira \rot{medido}. Coordenar sai quase tão caro quanto fazer.

\paragraph{Leitura.} \rot{inferido} É o que Kim et al.~\cite{kim2025} e a Vals.ai~\cite{vals2026}
preveem para uma tarefa que o agente sozinho já resolve: a equipe só acrescenta custo e tempo. A tarefa
do benchmark é pequena e sequencial (uma funcionalidade e duas correções no mesmo módulo). Ela não
mostra o que acontece numa tarefa que o agente sozinho não dá conta ou que se divide em partes
independentes. Esse é o caso que falta medir.

\section{O que isso muda no NarraTer}

\rot{inferido} A conclusão não é ``o NarraTer não deve orquestrar''. O transporte do NarraTer
(o terminal como API universal, a aresta como permissão, a resposta inteira lida do transcript, o
histórico gravado pelo backend) é exatamente o que a literatura pede de um orquestrador. A conclusão é
sobre a \emph{política}:

\begin{enumerate}
  \item \textbf{O padrão deve ser um agente só.} Montar uma equipe passa a exigir um motivo: duas ou
    mais subtarefas independentes com arquivos separados, um verificador com contexto limpo ou um
    contexto que não cabe num agente.
  \item \textbf{A equipe que mais vale em código é ``um escritor + um revisor''}, de preferência de
    fornecedores diferentes. É onde o NarraTer tem vantagem única: pôr um Codex revisando um Claude
    não depende de nenhum fornecedor.
  \item \textbf{``Pronto'' tem que ser verificado no ambiente}, rodando os testes ou o comando de
    aceite, e não declarado pelo agente.
  \item \textbf{O gerador de orquestração escolhe entre modelos prontos}, entre eles ``sozinho'', e
    registra por que montou ou não a equipe.
\end{enumerate}

\section{Limites deste documento}

\begin{itemize}
  \item A literatura foi lida, não reproduzida. Conferimos na fonte os números de Kim et al., da
    Vals.ai, da Cognition (2026) e do estudo de topologia; os demais estão citados com o link para
    conferir.
  \item Nossa medição usa uma tarefa só, já no teto de qualidade, e 3 rodadas por braço. Ela mede custo
    e tempo, não ganho de qualidade.
  \item A equipe usa Opus, Sonnet e Haiku; o agente sozinho usa só o Opus. Esta é a comparação que
    importa (``dou a tarefa ao Opus ou ao Opus com uma equipe?''), mas não isola o efeito de cada
    modelo.
  \item O braço da equipe é da versão G, e o braço sozinho é da versão atual (H). As duas diferem só
    no painel Activity, que não muda o protocolo.
\end{itemize}

\begin{thebibliography}{99}\small
\bibitem{kim2025} Y.~Kim et al. \emph{Towards a Science of Scaling Agent Systems}. arXiv:2512.08296, 2025.
  \url{https://arxiv.org/abs/2512.08296}
\bibitem{tran2026} D.~Tran e D.~Kiela. \emph{Single-Agent LLMs Outperform Multi-Agent Systems on Multi-Hop
  Reasoning Under Equal Thinking Token Budgets}. arXiv:2604.02460, 2026. \url{https://arxiv.org/abs/2604.02460}
\bibitem{rethinking2024} Q.~Wang et al. \emph{Rethinking the Bounds of LLM Reasoning: Are Multi-Agent
  Discussions the Key?} arXiv:2402.18272, 2024. \url{https://arxiv.org/abs/2402.18272}
\bibitem{mad2025} H.~Zhang et al. \emph{Stop Overvaluing Multi-Agent Debate: We Must Rethink Evaluation and
  Embrace Model Heterogeneity}. arXiv:2502.08788, 2025.
  \url{https://arxiv.org/abs/2502.08788}
\bibitem{vals2026} Vals.ai. \emph{Do Agent Teams Pay Off? A Case Study on Vibe Code Bench}, 2026-10-09.
  \url{https://www.vals.ai/blogs/multi-agent-vibe-code-bench}
\bibitem{agentless2024} C.~S.~Xia et al. \emph{Agentless: Demystifying LLM-based Software Engineering
  Agents}. arXiv:2407.01489, 2024. \url{https://arxiv.org/abs/2407.01489}
\bibitem{anthropic2025} Anthropic. \emph{How we built our multi-agent research system}, 2025.
  \url{https://www.anthropic.com/engineering/multi-agent-research-system}
\bibitem{topologia2026} \emph{An Empirical Study of Coordination Mode as the First-Class Citizen in
  From-Scratch Multi-Agent Coding}. arXiv:2607.27877, 2026. \url{https://arxiv.org/abs/2607.27877}
\bibitem{langchain2025} LangChain. \emph{Benchmarking Multi-Agent Architectures}, 2025.
  \url{https://www.langchain.com/blog/benchmarking-multi-agent-architectures}
\bibitem{cognition2025} W.~Yan (Cognition). \emph{Don't Build Multi-Agents}, 2025.
  \url{https://cognition.com/blog/dont-build-multi-agents}
\bibitem{cognition2026} W.~Yan (Cognition). \emph{Multi-Agents: What's Actually Working}, 2026-04-22.
  \url{https://cognition.com/blog/multi-agents-working}
\bibitem{mast2025} M.~Cemri et al. \emph{Why Do Multi-Agent LLM Systems Fail?} arXiv:2503.13657, 2025.
  \url{https://arxiv.org/abs/2503.13657}
\bibitem{falsosucesso2026} L.~Advani. \emph{From Confident Closing to Silent Failure: Characterizing False
  Success in LLM Agents}. arXiv:2606.09863, 2026. \url{https://arxiv.org/abs/2606.09863}
\bibitem{metagpt2023} S.~Hong et al. \emph{MetaGPT: Meta Programming for a Multi-Agent Collaborative
  Framework}. arXiv:2308.00352, 2023. \url{https://arxiv.org/abs/2308.00352}
\end{thebibliography}

\end{document}
